\documentclass[french]{llncs}

\usepackage[english, french]{babel}
\usepackage[utf8]{inputenc}

% --- Police ---
\usepackage{newtxtext}
\usepackage{newtxmath}
% ---------------------------------------------------------

\usepackage{graphicx}
\usepackage{amsmath}

\usepackage{listings}
\usepackage{subfig}
\usepackage{hyperref}
\usepackage{booktabs}

%-----------------------------------------
\title{Connaissant la latitude, la longitude et la population de sa ville, la structure topologique aide-t-elle à prédire le pays d'un aéroport~? \\
       Étude comparative sur un réseau aérien mondial}

\author{ Yennipena Jérémie \textsc{Coulibaly} \\
         Michel Junior \textsc{Delanoue} }

\institute{Université Claude Bernard Lyon 1, France \\
\url{https://github.com/cyjere/rendu-tp-gnn-Coulibaly-and-Delanoue}}

\lstset{
    basicstyle=\small,
    captionpos=b
}

\providecommand{\motcles}[1]{\textbf{Mots-clés:} #1}

\pagestyle{headings}
%-----------------------------------------

\begin{document}

    \renewcommand{\abstractname}{Abstract} \\
    \renewcommand{\thelstlisting}{\arabic{lstlisting}}

    \maketitle

    \begin{abstract}
        We study node classification on a global air transport network (3,363 airports,
        13,547 links) where the task is to predict the country of an airport from its
        attributes (longitude, latitude, population) and/or the graph of air connections.
        We compare classical machine learning and deep learning baselines
        (logistic regression, Random Forest, MLP) trained on node features alone with
        graph neural networks (GCN, GAT, SGC) trained either on structure alone or on
        features combined with structure, all evaluated with class-weighted losses and
        macro-F1 over five random splits to handle severe class imbalance (212 countries,
        ranging from 1 to 650 airports each). Contrary to the usual expectation, adding
        graph structure never improves the best feature-only model (Random Forest,
        macro-F1 $0.779 \pm 0.018$); structure alone performs no better than chance, and
        combining it with features systematically degrades the GNNs. We show that this is
        explained by low edge homophily (63.2\% of links connect two airports in the same
        country), concentrated on the largest international hubs, which dilutes the
        geographical signal on which message passing relies.

        \begin{keywords}
            graph neural networks, node classification, class imbalance,
            air transport network, ablation study
        \end{keywords}
    \end{abstract}

%-----------------------------------------

    \section{Introduction} \label{intro}

    Les réseaux aériens forment un graphe naturel~: les aéroports sont des nœuds, les liaisons
    directes des arêtes. Nous étudions si la \emph{structure} de ce graphe apporte une information
    utile pour prédire le \textbf{pays} d'un aéroport au-delà de ses seules caractéristiques
    individuelles, sur un graphe de 3\,363 aéroports et 13\,547 liaisons (longitude, latitude,
    population de la ville, pays cible).

    Notre problématique~: \emph{dans quelle mesure la topologie du réseau aérien améliore-t-elle
    la prédiction du pays par rapport aux seules caractéristiques individuelles~?} Nous comparons
    trois conditions~: (A) modèles classiques ML/DL sur les seules caractéristiques~; (B) GNN
    entraînés uniquement sur la structure~; (C) mêmes GNN sur caractéristiques \emph{et} structure.
    Cette décomposition isole la contribution propre de la topologie.

    \section{État de l'art} \label{sect:eda}

    Pour la condition graphe, nous retenons trois architectures. Le \textbf{GCN} \cite{Kipf2017}
    moyenne la représentation d'un nœud avec celle de ses voisins, pondérée par le degré. Le
    \textbf{SGC} \cite{Wu2019} supprime les non-linéarités et matrices de poids intermédiaires
    pour ne garder que la propagation, ce qui teste l'apport réel de la non-linéarité. Le
    \textbf{GAT} \cite{Velickovic2018} remplace la moyenne par un poids d'attention appris par
    paire de nœuds~: nous testons s'il peut ignorer des voisins non pertinents en cas de faible
    homophilie.

    Pour les baselines sans graphe, le \textbf{Random Forest} \cite{Breiman2001} est un
    classifieur non linéaire robuste adapté aux petites features continues, et le
    \textbf{perceptron multicouche} \cite{Rumelhart1986} sert de baseline DL sans structure.
    Le déséquilibre sévère (212 pays, 1 à 650 aéroports) justifie, d'après \cite{Wang2021},
    une pondération de la perte par classe et une évaluation en macro-F1 plutôt qu'en accuracy.

    \section{Données et protocole} \label{sect:protocol}

    \paragraph{Données.} Le graphe est chargé depuis un fichier GraphML~: 3\,363 nœuds (aéroports),
    13\,547 arêtes (liaisons), avec pour chaque nœud \texttt{lon}, \texttt{lat}, \texttt{population},
    \texttt{city\_name} et \texttt{country} (cible). La cible comporte 212 classes brutes très
    déséquilibrées (USA~: 650 aéroports, médiane~: 4, 62 pays à un seul aéroport), ce qui rend
    impossible toute évaluation sur les plus petites classes.

    \paragraph{Filtrage et split.} Nous retirons les pays à moins de 3 aéroports
    (\texttt{MIN\_CLASS\_SIZE=3}), ramenant le problème à 126 classes sur 3\,253 nœuds et
    12\,272 arêtes. Un split stratifié classique échoue encore sur ce seuil~: après une première
    coupe train/reste, de nombreuses classes à 3 éléments n'ont plus qu'un seul exemple,
    impossible à répartir entre validation et test. Nous utilisons donc un split par classe
    garantissant au moins un exemple dans chacun des trois ensembles, avec un repli
    proportionnel pour les classes plus grandes~: 69,0\,\%~/~15,5\,\%~/~15,5\,\%, répété sur 5 seeds.

    \paragraph{Features.} Les coordonnées brutes \texttt{lon/lat} posent deux problèmes~:
    longitude circulaire et singularités aux pôles. Nous testons donc un encodage sur la sphère
    unité $(x,y,z) = (\cos(\text{lat})\cos(\text{lon}), \cos(\text{lat})\sin(\text{lon}), \sin(\text{lat}))$,
    combiné à $\log(1+\text{population})$, standardisé sur le train uniquement. Cet encodage
    améliore tous les modèles (section~\ref{sect:results}) et devient la représentation par défaut.

    \paragraph{Déséquilibre et hyperparamètres.} Les modèles sklearn utilisent
    \texttt{class\_weight="balanced"}~; les réseaux de neurones (MLP, GCN, GAT, SGC) une
    \texttt{CrossEntropyLoss} pondérée par l'inverse de la fréquence de classe (calculée sur le
    train de chaque seed). L'architecture du MLP (256 unités, 2 couches, \texttt{lr=0,005},
    \texttt{dropout=0,1}) a été choisie par recherche sur la validation du seed~0 uniquement
    (macro-F1~: 0,593~$\rightarrow$~0,684), jamais sur le test.

    \section{Expériences et résultats} \label{sect:results}

    Toutes les métriques sont moyennées sur 5 seeds (moyenne~$\pm$~écart-type), macro-F1 étant
    la métrique principale.

    \subsection{Condition A~: caractéristiques seules}

    Le tableau~\ref{tab:features} compare les modèles sans graphe. L'encodage \texttt{xyz}
    améliore tous les modèles par rapport à \texttt{lon/lat} (RF pondéré~: $0,577 \rightarrow 0,779$),
    probablement parce qu'il évite les discontinuités aux frontières et capture mieux les distances
    relatives. La pondération de la perte aide nettement le Random Forest ($+0,04$ de macro-F1)
    mais pas le MLP, dont les deux versions sont équivalentes aux écarts-types près~: son effet
    dépend donc du modèle et ne doit pas être supposé systématique. Ce résultat nuance l'idée reçue
    selon laquelle la pondération est toujours bénéfique.

    \begin{table}[ht]
     \begin{center}
       \scriptsize
       \setlength{\tabcolsep}{4pt}
       \renewcommand{\arraystretch}{0.8}
       \begin{tabular}{l c c c}
       \toprule
       Modèle (features) & macro-F1 & balanced acc. & accuracy \\
       \midrule
       LogReg (xyz)                       & $0,285 \pm 0,011$ & $0,358 \pm 0,008$ & $0,491 \pm 0,013$ \\
       RF pondéré (lon/lat)               & $0,577 \pm 0,007$ & $0,631 \pm 0,009$ & $0,771 \pm 0,010$ \\
       \textbf{RF pondéré (xyz)}          & $\mathbf{0,779 \pm 0,018}$ & $\mathbf{0,816 \pm 0,013}$ & $0,868 \pm 0,010$ \\
       RF non pondéré (xyz)               & $0,741 \pm 0,026$ & $0,769 \pm 0,025$ & $\mathbf{0,886 \pm 0,011}$ \\
       MLP pondéré (xyz)                  & $0,613 \pm 0,011$ & $0,672 \pm 0,014$ & $0,785 \pm 0,009$ \\
       MLP non pondéré (xyz)              & $0,618 \pm 0,035$ & $0,658 \pm 0,035$ & $0,846 \pm 0,019$ \\
       \bottomrule
       \end{tabular}
      \caption{Condition A (features seules), 5 seeds. Meilleur macro-F1 en gras.} \label{tab:features}
     \end{center}
    \end{table}

    \subsection{Conditions B et C~: structure seule, puis features + structure}

    Le tableau~\ref{tab:graph} ajoute les GNN, entraînés soit sur une feature constante
    (B~: topologie seule), soit sur les features \texttt{xyz} (C~: features + structure), avec
    la même perte pondérée et le même protocole à 5 seeds.

    \begin{table}[ht]
     \begin{center}
       \scriptsize
       \setlength{\tabcolsep}{4pt}
       \renewcommand{\arraystretch}{0.8}
       \begin{tabular}{l c c c}
       \toprule
       Modèle (condition) & macro-F1 & balanced acc. & accuracy \\
       \midrule
       \textbf{RF pondéré (A, xyz)} & $\mathbf{0,779 \pm 0,018}$ & $\mathbf{0,816 \pm 0,013}$ & $\mathbf{0,868 \pm 0,010}$ \\
       GCN (B, topologie seule)     & $0,003 \pm 0,004$ & $0,013 \pm 0,009$ & $0,007 \pm 0,005$ \\
       GAT (B, topologie seule)     & $0,000 \pm 0,000$ & $0,008 \pm 0,000$ & $0,012 \pm 0,006$ \\
       SGC (B, topologie seule)     & $0,000 \pm 0,000$ & $0,006 \pm 0,003$ & $0,002 \pm 0,001$ \\
       GCN (C, features+structure)  & $0,493 \pm 0,020$ & $0,552 \pm 0,023$ & $0,698 \pm 0,011$ \\
       GAT (C, features+structure)  & $0,252 \pm 0,042$ & $0,331 \pm 0,038$ & $0,474 \pm 0,049$ \\
       SGC (C, features+structure)  & $0,370 \pm 0,031$ & $0,426 \pm 0,032$ & $0,579 \pm 0,032$ \\
       \bottomrule
       \end{tabular}
      \caption{Conditions B et C, 5 seeds, mêmes features \texttt{xyz} que le Random Forest de référence (A).} \label{tab:graph}
     \end{center}
    \end{table}

    Deux constats. La topologie seule (B) ne porte quasiment aucun signal~: les macro-F1
    ($0,000$ à $0,003$) sont au niveau du hasard sur 126 classes ($1/126 \approx 0,008$).
    La structure du réseau aérien, prise isolément, n'apporte donc aucune information
    discriminante~: les connexions aériennes ne sont pas un proxy suffisant de l'appartenance
    nationale. Ajouter la structure aux features (C) ne bat jamais le Random Forest de la
    condition A et dégrade même les performances par rapport aux mêmes features sans graphe~:
    le meilleur GNN (GCN) perd 0,29 point de macro-F1. Ce phénomène s'explique par la
    propagation de bruit~: les voisins étrangers injectent des informations contradictoires
    qui perturbent la représentation apprise, alors que le Random Forest, ignorant le graphe,
    se concentre sur les seules caractéristiques géographiques.

    \section{Discussion et conclusion} \label{sect:discussion}

    \paragraph{Pourquoi la structure n'aide pas.} Les GNN à passage de messages supposent
    implicitement que les voisins partagent la même classe. Or seulement \textbf{63,2\,\%} des
    arêtes (7\,761 sur 12\,272) relient deux aéroports du même pays. Ce déficit d'homophilie est
    concentré sur les plus gros hubs, qui pèsent le plus dans l'agrégation~: Paris (80\,\% de
    voisins étrangers), Londres (93\,\%), Francfort (91\,\%), Amsterdam (99\,\%) — contre 24\,\%
    pour Chicago, à dominante domestique. Les GNN sont donc pénalisés par les hubs fortement
    connectés à l'international, qui brouillent le signal local~; à l'inverse, un pays à réseau
    domestique dense comme les États-Unis pourrait bénéficier d'une agrégation plus homogène.

    \paragraph{Diagnostic par classe.} Les erreurs du meilleur modèle (Random Forest, xyz) ne sont
    pas aléatoires~: le macro-F1 moyen passe de $0,674$ pour les classes à $\leq 5$ aéroports
    d'entraînement (n=46) à $0,858$ au-delà (n=80), et 7 des 8 pires classes sont confondues avec
    leur voisin géographique réel et plus représenté (Bosnie~$\rightarrow$~Croatie,
    Cambodge~$\rightarrow$~Thaïlande, Cameroun~$\rightarrow$~Nigeria). Un petit pays n'échoue donc
    pas par manque de données seul, mais quand sa frontière de décision est absorbée par un grand
    voisin dans l'espace des features.

    \paragraph{Conclusion.} Sur ce jeu de données, la structure topologique du réseau aérien
    \textbf{n'améliore pas} la prédiction du pays~: elle la dégrade, car la connectivité aérienne
    est largement internationale et la position géographique porte déjà l'essentiel du signal.
    Limites~: un seul graphe/une seule tâche, d'autres architectures non testées (GraphSAGE, APPNP),
    et l'observation sur Chicago à vérifier systématiquement.

    \newpage
    \bibliographystyle{splncs03}
    \bibliography{bibM2.bib}

\end{document}